% from main.tex
%%%%%%%%%%  scan idx 282  |  folio 269  %%%%%%%%%%
compatibilité du rang avec les images inverses (I 6.12).

(ii) Quitte à tronquer et à translater, on peut, pour la première assertion, supposer \(E\) strictement parfait, \(p=0\), et \(E^i=0\) pour \(i\notin[-1,+1]\). Pour \(x\in\pt(X)\), posons \(E_x\otimes_{\mathcal O_{X,x}}^{\mathrm L}k(x)=E_{k(x)}\). On a d'après (5.5.1)
\[
\begin{gathered}
 \operatorname{rg}_{k(x)}H^0(E_{k(x)})=\operatorname{rg}_{k(x)}H^{-1}(E_{k(x)})+\operatorname{rg}_{k(x)}H^1(E_{k(x)})\\
 +\operatorname{rg}(E_x),
\end{gathered}
\]
donc il suffit de prouver que la fonction \(x\mapsto\operatorname{rg}_{k(x)}H^{-1}(E_{k(x)})\) (resp. \(\operatorname{rg}_{k(x)}H^1(E_{k(x)})\)) est semi-continue supérieurement (puisque la fonction \(x\mapsto\operatorname{rg}(E_x)=\operatorname{rg}(E_{k(x)})\) est localement constante). Pour la fonction \(x\mapsto\operatorname{rg}_{k(x)}H^1(E_{k(x)})\), c'est une conséquence immédiate de (5.4.1 b))\footnote{compte tenu de ce que l'on a \(H^1(E_{k(x)})\simeq H^1(E)_{k(x)}\).}. Pour la fonction \(x\mapsto\operatorname{rg}_{k(x)}H^{-1}(E_{k(x)})\), on peut supposer \(E^1=0\), i.e. remplacer \(E\) par \(E'=(0\to E^{-1}\to E^0\to0)\). On sait déjà que la fonction \(x\mapsto\operatorname{rg}_{k(x)}H^0(E'_{k(x)})\) est semi-continue supérieurement, donc il en est de même de la fonction \(x\mapsto\operatorname{rg}_{k(x)}H^{-1}(E'_{k(x)})=\operatorname{rg}_{k(x)}H^{-1}(E_{k(x)})\), puisque la fonction \(x\mapsto\operatorname{rg}_{k(x)}H^0(E'_{k(x)})-\operatorname{rg}_{k(x)}H^{-1}(E'_{k(x)})=\operatorname{rg}(E'_x)\) est localement constante. Prouvons maintenant la seconde assertion. Quitte à tronquer, on peut supposer \(E\) strictement parfait et \(E^i=0\) pour \(i<p-1\). Le \(\mathcal O_X\)-Module \(Z'^i(E)=\Coker(E^{i-1}\to E^i)\) est donc de présentation finie. On a une suite exacte :
\begin{equation}
 0\to H^p(E)\to Z'^p(E)\to E^{p+1}\to Z'^{p+1}(E)\to0.\tag{$*$}
\end{equation}
Comme le foncteur \(M\mapsto H^p(E\otimes^{\mathrm L}M)\) est exact à droite et à gauche, \(Z'^{p+1}(E)\) et \(Z'^p(E)\)) sont plats d'après (5.2 (iii)), donc (I 5.8.3) localement libre de type fini. L'exactitude de (*) implique alors que \(H^p(E)\) est localement libre de type fini, et par suite que la fonction \(x\mapsto\operatorname{rg}_{k(x)}H^p(E_{k(x)})=\operatorname{rg}(H^p(E)_x)\) est
